\subsection{The Lord Is Nigh}\label{sec:third-week}

The Mass of the Third Sunday begins with one word, \latin{Gaudete},
``Rejoice'', and the Epistle supplies the verse in which the word stands, and
the reason for it:

\begin{quote}
Rejoice in the Lord always: again, I say, rejoice. Let your modesty be known to
all men. The Lord is nigh.\par
{\smallskip\noindent\small ---Philippians 4:4--5\par}
\end{quote}

\latin{Dominus prope est}: the Lord is near. A Christian who has been fasting
and watching for two weeks is told why: the One he waits for is close.

The church shows it. On this Sunday, in its Office and Mass only, the priest may
wear rose, and the organ, silent through the rest of Advent except at
Benediction or to sustain the singing, may be played, though no other
instrument. The station is St Peter's. The Sundays of Advent are of the
first class; the weekdays are of the third, except the Ember Days of this week
and the weekdays from 17 December, which are of the second.

\subsubsection*{Who art thou?}

The Gospel is a cross-examination. Priests and Levites come from Jerusalem to ask
John who he is. He tells them what he is not. ``I am not the
Christ.'' He is not Elias; he is not the prophet. They press him for something
to take back to those who sent them, and he hands them a verse of Isaiah:

\begin{quote}
He said: I am the voice of one crying in the wilderness, make straight the way
of the Lord, as said the prophet Isaias.\par
{\smallskip\noindent\small ---John 1:23\par}
\end{quote}

St Gregory the Great says that John calls himself the voice because a voice
sounds first and the word is heard after it: John goes before the Word. In his
homily on this Gospel Gregory begins with John's lowliness:

\begin{quote}
Dearly beloved brethren, the first thing which striketh us in to-day's Gospel is
the lowly-mindedness of John. He was so great that it was thought he might be
the Christ; yet he soberly chose rather to seem only what he really was, than to
let the belief of men invest him with a dignity which did not belong to him; for
``he confessed, and denied not, but confessed, I am not the Christ,'' at the same
time he would not deny what he was in reality; and thus his very truth-speaking
made him a member of Him Whose title he would not by falsehood take.\par
{\smallskip\noindent\small ---St Gregory the Great, Homily 7 on the Gospels\par}
\end{quote}

Then John says to the men from Jerusalem: ``there hath stood one in the midst of
you, whom you know not.'' On the first night of Advent the Church read the
prophet's complaint, ``Israel hath not known me''. Now the One whom Israel did
not know is standing among them, unrecognised, and the voice in the desert is
there to point him out.

St Augustine, preaching on the same chapter, saw how much depended on John's
humility, and how much on Christ's:

\begin{quote}
You have very often heard, holy brethren, and you know well, that John the
Baptist, in proportion as he was greater than those born of women, and was more
humble in his acknowledgment of the Lord, obtained the grace of being the
friend of the Bridegroom; zealous for the Bridegroom, not for himself; not
seeking his own honor, but that of his Judge, whom as a herald he preceded.
\dots\ For He had come humbly and concealed from the first; the more concealed
in proportion as He was more humble: but the people, despising in their pride
the humility of God, crucified their Saviour, and made Him their condemner.\par
{\smallskip\noindent\small ---St Augustine, Tractate 4 on St John\par}
\end{quote}

The Lord is near, and he is hidden, and this week the Church prays for light:

\begin{studyframe}
\latin{Aurem tuam, quaesumus, Domine, precibus nostris accommoda: et mentis
nostrae tenebras gratia tuae visitationis illustra: Qui vivis et regnas cum Deo
Patre in unitate.}
\medskip

Incline Thine ear to our prayers, O Lord, we beseech Thee; and make bright the
darkness of our minds by the grace of Thy visitation. Who livest.\par
\end{studyframe}

The week of the Ember Days is the Church's old week of penance, and the week to
go to confession. Before the Ember Days begin on Wednesday, put John's
questions to yourself. Who am I? What do I claim to be that I am not? What have
I denied that I am? Then take up the examination of conscience in
Section~\ref{sec:confession}, and go.
